% 									Resumo em português

\setlength{\absparsep}{18pt} % ajusta o espaçamento dos parágrafos do resumo
\begin{resumo}
	A utilização de simulações computacionais tornou-se uma etapa fundamental no desenvolvimento e na análise de experimentos radioastronômicos modernos, permitindo investigar o impacto de efeitos instrumentais, testar algoritmos de processamento de dados e avaliar estratégias observacionais em ambientes controlados. Nesse contexto, esta dissertação apresenta o desenvolvimento de uma pipeline de simulação voltada para observações radioastronômicas do tipo \textit{Single Dish}, com foco na geração de \textit{Time Ordered Data} (TOD) sintéticos e na reconstrução de mapas do céu por meio de técnicas de \textit{mapmaking}.
	A pipeline desenvolvida modela as principais etapas da cadeia observacional, incluindo a representação do céu astrofísico, a resposta instrumental, a geração de dados temporais e a reconstrução de mapas. Como aplicação, o sistema foi concebido para reproduzir características relevantes do radiotelescópio BINGO (\textit{Baryon Acoustic Oscillations from Integrated Neutral Gas Observations}), experimento dedicado a observações da emissão integrada de hidrogênio neutro utilizando a técnica de \textit{Intensity Mapping}.
	A arquitetura do software foi desenvolvida com base em princípios de Programação Orientada a Objetos e \textit{Domain Driven Design}, visando promover modularidade, extensibilidade e facilidade de manutenção. A validação dos módulos implementados foi realizada por meio de cenários de teste controlados e comparações com modelos teóricos e simulações independentes.
	Os resultados obtidos indicam que a pipeline é capaz de reproduzir aspectos fundamentais de observações radioastronômicas \textit{Single Dish}, incluindo a geração consistente de \textit{Time Ordered Data}, a incorporação de efeitos instrumentais e a reconstrução de mapas do céu. Dessa forma, o trabalho disponibiliza uma infraestrutura computacional reutilizável para estudos de instrumentação, desenvolvimento de metodologias de processamento de dados e suporte a atividades de simulação associadas ao projeto BINGO e a experimentos observacionais semelhantes.


\vspace{\onelineskip}
	\noindent
	\textbf{Palavras-chave}: Radioastronomia, Single Dish, Time Ordered Data, Mapmaking, Simulação, Pipeline.
\end{resumo}

% 										Resumo em inglês

\begin{resumo}[Abstract]
\begin{otherlanguage*}{english}
	The use of computational simulations has become a fundamental component in the development and analysis of modern radio astronomy experiments, enabling the investigation of instrumental effects, the testing of data processing algorithms, and the evaluation of observational strategies within controlled environments. In this context, this dissertation presents the development of a simulation pipeline for Single-Dish radio astronomy observations, with emphasis on the generation of synthetic Time Ordered Data (TOD) and the reconstruction of sky maps through mapmaking techniques.
The developed pipeline models the main stages of the observational chain, including the representation of the astrophysical sky, the instrumental response, the generation of time-domain data, and map reconstruction. As a practical application, the system was designed to reproduce relevant characteristics of the BINGO (Baryon Acoustic Oscillations from Integrated Neutral Gas Observations) radio telescope, an experiment dedicated to observations of the integrated neutral hydrogen emission using the Intensity Mapping technique.
The software architecture was developed based on Object-Oriented Programming and Domain-Driven Design principles, aiming to promote modularity, extensibility, and maintainability. The implemented modules were validated through controlled test scenarios and comparisons with theoretical models and independent simulations.
The obtained results indicate that the pipeline is capable of reproducing fundamental aspects of Single-Dish radio astronomy observations, including the consistent generation of Time Ordered Data, the incorporation of instrumental effects, and the reconstruction of sky maps. Therefore, this work provides a reusable computational infrastructure for instrumentation studies, the development of data processing methodologies, and simulation activities related to the BINGO project and similar observational experiments.

	
\vspace{\onelineskip}
	\noindent 
	\textbf{Keywords}: Radio Astronomy, Single-Dish, Time Ordered Data, Mapmaking, Simulation, Pipeline.
\end{otherlanguage*}
\end{resumo}